% part2_formal_results.tex
% Part II of the online appendix. Includable file: no preamble.
% Requires part2_macros.tex in the preamble of the host document, natbib,
% booktabs, and xr-hyper with \externaldocument[P-]{<path to the paper>/main}.
% Labels defined here carry the prefix "fr:". Labels with the prefix "P-"
% are the paper's; labels with the prefixes "oa:" and "oatab:" are defined
% in Part I of the host document.
% Every printed number is asserted by build/check_part2_numbers.py.

\part*{Part II. Formal results}
\addcontentsline{toc}{part}{Part II. Formal results}

\noindent This part states and proves the results that
Sections~\ref{P-sec:process}, \ref{P-sec:measurement},
and~\ref{P-sec:mediation} of the paper use without proof.
Section~\ref{fr:sec:atten} derives the attenuation schedule of
Section~\ref{P-sec:measurement} of the paper from the estimated earnings
process: the exact adjustment for serial correlation, the number of waves
needed to reach a target share of the coefficient, and a closed form for
the simulated rank-rank factors. Section~\ref{fr:sec:md} treats the
minimum-distance estimator of Section~\ref{P-sec:md} of the paper:
identification, the ridge of Figure~\ref{P-fig:ridge} of the paper, the
distribution of the criterion under fixed weights, and a decomposition of
the criterion into a stationarity component and a shape component.
Section~\ref{fr:sec:identity} writes the accounting identity of
Section~\ref{P-sec:mediation} of the paper in potential outcomes and shows
what its remainder contains. Every number that is not copied from the
paper is produced by \code{formal_results/code/appendix_theory_computations.py} in the
repository, and \code{online_appendix/build/check_part2_numbers.py} asserts each
printed number against its source.

\section{The attenuation schedule}
\label{fr:sec:atten}

\subsection{Setup}

Residual log earnings of child $i$ in wave $t$ follow
equation~\eqref{P-eq:feAR} of the paper,
\begin{equation}
\tilde y_{it} = \eta_i + v_{it}, \qquad
v_{it} = \rho\, v_{i,t-1} + \varepsilon_{it}, \qquad |\rho| < 1,
\label{fr:eq:model}
\end{equation}
with $\eta_i$ a permanent component, $\varepsilon_{it}$ serially
uncorrelated with mean zero and variance $\sigma^2_\varepsilon$, and
$\{\varepsilon_{it}\}$ independent of $\eta_i$. The persistence parameter
is written $\rho_v$ in the paper; this part writes $\rho$. Under
stationarity the transitory component has variance
$\sigma^2_v = \sigma^2_\varepsilon/(1-\rho^2)$ and autocovariance function
$\gamma_k = \rho^{k}\sigma^2_v$, $k \ge 0$, so the autocovariances of
$\tilde y$ are those of equation~\eqref{P-eq:moments} of the paper,
\begin{equation}
\omega_k \;\equiv\; \Cov(\tilde y_{it}, \tilde y_{i,t+k}) \;=\; \sigma^2_\eta + \rho^{k}\,\sigma^2_v .
\label{fr:eq:omega}
\end{equation}
One period is one CFPS wave (two years). The permanent share is
$s = \sigma^2_\eta/(\sigma^2_\eta + \sigma^2_v)$. Table~\ref{P-tab:md} of
the paper estimates $(\hat\sigma^2_\eta, \hat\sigma^2_\varepsilon, \hat\rho) =
(0.125, 0.326, 0.16)$ on the pooled sample, so $\hat\sigma^2_v = 0.335$ and
$\hat s = 0.27$.

\subsection{The variance of a wave average}

\begin{proposition}[Variance of a $T$-wave mean of a stationary AR(1)]
\label{fr:prop:fT}
Let $\bar v_{iT} = T^{-1}\sum_{t=1}^{T} v_{it}$. Then
\begin{equation}
\Var(\bar v_{iT}) \;=\; \frac{\sigma^2_v}{T}\, f_T(\rho), \qquad
f_T(\rho) \;=\; 1 + 2\sum_{k=1}^{T-1}\Big(1 - \frac{k}{T}\Big)\rho^{k}
\;=\; \frac{1+\rho}{1-\rho} \;-\; \frac{2\rho\,(1-\rho^{T})}{T\,(1-\rho)^{2}} .
\label{fr:eq:fT}
\end{equation}
For $\rho \in (0,1)$: $f_1 = 1$; $f_T$ is strictly increasing in $T$ with
limit $(1+\rho)/(1-\rho)$; $f_T/T$ is strictly decreasing in $T$; and for
$T \ge 2$, $f_T$ is strictly increasing in $\rho$.
\end{proposition}

\begin{proof}
$\Var(\bar v_{iT}) = T^{-2}\sum_{s=1}^{T}\sum_{t=1}^{T}\gamma_{|s-t|}
= T^{-2}\big[T\gamma_0 + 2\sum_{k=1}^{T-1}(T-k)\gamma_k\big]$, and
substituting $\gamma_k = \rho^k\sigma^2_v$ gives the sum form of $f_T$.
For the closed form, use $\sum_{k=1}^{T-1}\rho^k = \rho(1-\rho^{T-1})/(1-\rho)$
and $\sum_{k=1}^{T-1}k\rho^k = \rho\big(1 - T\rho^{T-1} + (T-1)\rho^{T}\big)/(1-\rho)^2$;
collecting terms,
\[
f_T = 1 + \frac{2\rho(1-\rho^{T-1})}{1-\rho} - \frac{2\rho\big(1 - T\rho^{T-1} + (T-1)\rho^{T}\big)}{T(1-\rho)^2},
\]
and the two fractions combine to $\frac{2\rho}{1-\rho} - \frac{2\rho(1-\rho^T)}{T(1-\rho)^2}$,
which with the leading 1 is \eqref{fr:eq:fT}. In the sum form each summand
$(1-k/T)\rho^k$ is positive and increasing in $T$ and in $\rho$, and the
number of summands grows with $T$, which gives the two monotonicity
statements for $f_T$; the limit follows from the closed form. For
$f_T/T$, write the mean of $T+1$ consecutive terms as the average of its
$T+1$ leave-one-out means. Each leave-one-out mean averages $T$ terms
whose pairwise distances in time are at least those of $T$ consecutive
terms, so with $\rho \in (0,1)$ its variance is at most
$\Var(\bar v_{iT})$. The variance of an average of random variables is at
most the largest of their variances, strictly so unless they are
perfectly correlated, which the leave-one-out means are not. Hence
$\Var(\bar v_{i,T+1}) < \Var(\bar v_{iT})$.
\end{proof}

At the pooled estimate $\hat\rho = 0.16$, $f_2 = 1.16$, $f_5 = 1.29$,
$f_9 = 1.33$, and $f_\infty = 1.38$. Because transitory components of consecutive waves are correlated, averaging removes their variance more slowly than the independence calculation implies; the variance of a long average is up to 38 percent larger.

\subsection{Attenuation}

\begin{assumption}[Classical measurement structure]
\label{fr:ass:cev}
The outcome satisfies $c_i = a + b\,y^*_i + e_i$, where $y^*_i$ is
permanent income with variance $\sigma^2_\eta$ and $b$ is the coefficient
of the linear projection of $c_i$ on $y^*_i$, so that
$\Cov(y^*_i, e_i) = 0$. The researcher observes
$x_{iT} = y^*_i + \bar v_{iT}$, with $\bar v_{iT}$ as in
Proposition~\ref{fr:prop:fT}, and regresses $c_i$ on $x_{iT}$ by least
squares on an i.i.d.\ sample. The transitory components are uncorrelated
with $y^*_i$ and with $e_i$.
\end{assumption}

\begin{proposition}[The attenuation schedule]
\label{fr:prop:eta}
Under \eqref{fr:eq:model} and Assumption~\ref{fr:ass:cev}, the least
squares slope $\hat b_T$ satisfies
\begin{equation}
\plim\,\hat b_T \;=\; b\,\lambda_T, \qquad
\lambda_T \;=\; \frac{\sigma^2_\eta}{\sigma^2_\eta + \sigma^2_v\, f_T(\rho)/T}.
\label{fr:eq:eta}
\end{equation}
The independence approximation of equation~\eqref{P-eq:eta} of the paper
is $\lambda^{\circ}_T = \sigma^2_\eta/(\sigma^2_\eta + \sigma^2_v/T)$, that is,
\eqref{fr:eq:eta} with $f_T$ replaced by one. For $\rho \in (0,1)$ and
given $(\sigma^2_\eta, \sigma^2_v)$:
(i) $\lambda_T \le \lambda^{\circ}_T$, with equality only at $T = 1$;
(ii) $\lambda_T$ is strictly increasing in $T$ with limit one;
(iii) $1 - \lambda_T = O(1/T)$; and
(iv) for $T \ge 2$, $\lambda_T$ is strictly decreasing in $\rho$.
\end{proposition}

\begin{proof}
The least squares slope converges to $\Cov(c, x_T)/\Var(x_T) =
b\,\sigma^2_\eta/(\sigma^2_\eta + \Var(\bar v_T))$ because
$\Cov(y^*, e) = \Cov(y^*, \bar v_T) = \Cov(e, \bar v_T) = 0$;
substitute Proposition~\ref{fr:prop:fT}. (i) follows from $f_T \ge 1$ with
equality only at $T = 1$. (ii) and (iii) follow from $f_T/T$ decreasing
in $T$ and bounded above by $f_\infty/T$. (iv) follows from $f_T$
increasing in $\rho$.
\end{proof}

\begin{table}[!htb]
\centering
\small
\caption{The attenuation schedule with bootstrap intervals.}
\label{fr:tab:etaboot}
\begin{tabular}{@{}cccccc@{}}
\toprule
& \multicolumn{2}{c}{Independence approximation $\lambda^{\circ}_T$} & \multicolumn{2}{c}{Exact $\lambda_T$} & \\
\cmidrule(lr){2-3}\cmidrule(lr){4-5}
Waves $T$ & Point & 95 percent interval & Point & 95 percent interval & $1/\lambda_T$ \\
\midrule
1 & 0.271 & [0.23, 0.31] & 0.271 & [0.23, 0.31] & 3.69 \\
2 & 0.427 & [0.37, 0.47] & 0.391 & [0.33, 0.45] & 2.56 \\
3 & 0.528 & [0.47, 0.57] & 0.476 & [0.41, 0.54] & 2.10 \\
4 & 0.598 & [0.54, 0.64] & 0.540 & [0.47, 0.60] & 1.85 \\
5 & 0.651 & [0.60, 0.69] & 0.590 & [0.52, 0.65] & 1.69 \\
6 & 0.691 & [0.64, 0.73] & 0.631 & [0.56, 0.69] & 1.58 \\
7 & 0.723 & [0.68, 0.76] & 0.664 & [0.59, 0.72] & 1.51 \\
9 & 0.770 & [0.73, 0.80] & 0.716 & [0.65, 0.77] & 1.40 \\
\bottomrule
\end{tabular}
\par\smallskip
\begin{minipage}{0.92\textwidth}
\footnotesize\emph{Notes:} Equation \eqref{fr:eq:eta} at the pooled
minimum-distance estimate of Table~\ref{P-tab:md} of the paper. Intervals
are percentile intervals from 1,000 bootstrap replications that resample
children, re-demean the panel, recompute the fifteen pairwise
autocovariances, and re-solve the minimum-distance fit in every draw, then
evaluate \eqref{fr:eq:eta} at each replication's parameters. The intervals
for the exact factor are wider because $f_T(\rho)$ adds the sampling
variation of $\hat\rho$ to that of the two variances. The last column is
the factor by which a researcher would multiply an estimated coefficient
to undo attenuation under the exact formula, were the parameters known.
\end{minipage}
\end{table}

Table~\ref{fr:tab:etaboot} attaches sampling uncertainty to the schedule
of Table~\ref{P-tab:eta} of the paper. The single-wave factor is 0.27 with a 95 percent interval of $[0.23, 0.31]$: even at the top of the interval a single wave recovers less than a third of the coefficient on permanent income. At seven waves the exact factor is 0.66 with interval $[0.59, 0.72]$. The paper's statement that six biennial waves reach two thirds uses the independence approximation; under the exact formula the crossing is at eight waves, as the footnote to
equation~\eqref{P-eq:eta} of the paper states.

\subsection{Which side of the regression}
\label{fr:sec:side}

Suppose the outcome is also observed as a $T_c$-wave average,
$c_{iT_c} = c_i + \bar w_{iT_c}$, where $\bar w_{iT_c}$ is uncorrelated
with $y^*_i$, $\bar v_{iT}$, and $e_i$. The slope of $c_{iT_c}$ on
$x_{iT}$ then has probability limit $b\,\lambda_{T}$ regardless of $T_c$:
noise in the dependent variable adds to the regression disturbance and is
uncorrelated with the regressor, which leaves the probability limit of
the slope unchanged, whereas noise in the regressor enters the denominator
as in Proposition~\ref{fr:prop:eta}. If the roles are reversed and $x_{iT}$
is regressed on $c_{iT_c}$, the slope is the coefficient of the projection
of $y^*_i$ on $c_i$ multiplied by $\Var(c_i)/\Var(c_{iT_c})$, the
reliability of the averaged outcome, and $\lambda_T$ does not enter.

Section~\ref{P-sec:measurement} of the paper accordingly takes the
relevant $T$ in an intergenerational elasticity regression to be the
number of parental income observations, and panel B of Table~\ref{P-tab:eta} of
the paper recomputes the correction at the earnings process of linked
parents, whose permanent share is 0.43 against the children's 0.27
(Table~\ref{oatab:parent}).

\subsection{Waves needed to reach a target}

\begin{proposition}[Waves required]
\label{fr:prop:tstar}
For a target share $\kappa \in (0,1)$ of the coefficient define
$T^*(\kappa) = \min\{T : \lambda_T \ge \kappa\}$ and
$T^{\circ *}(\kappa) = \min\{T : \lambda^{\circ}_T \ge \kappa\}$. Then
\begin{equation}
T^{\circ *}(\kappa) \;=\; \Big\lceil \frac{\kappa}{1-\kappa}\cdot\frac{\sigma^2_v}{\sigma^2_\eta} \Big\rceil
\;=\; \Big\lceil \frac{\kappa}{1-\kappa}\cdot\frac{1-s}{s} \Big\rceil ,
\label{fr:eq:tstar}
\end{equation}
and for $\rho \in (0,1)$, $T^*(\kappa) \ge T^{\circ *}(\kappa)$, with
$T^*(\kappa)$ the smallest $T$ such that
$f_T(\rho)/T \le (1-\kappa)\,\sigma^2_\eta/(\kappa\,\sigma^2_v)$.
\end{proposition}

\begin{proof}
$\lambda^{\circ}_T \ge \kappa$ is equivalent to
$\sigma^2_\eta(1-\kappa) \ge \kappa\,\sigma^2_v/T$, that is
$T \ge \kappa\sigma^2_v/((1-\kappa)\sigma^2_\eta)$; the second expression
substitutes $\sigma^2_v/\sigma^2_\eta = (1-s)/s$. The exact statement
rearranges \eqref{fr:eq:eta} in the same way with $f_T/T$ in place of
$1/T$, and $f_T \ge 1$ gives the inequality between the two thresholds.
\end{proof}

At the pooled estimates and $\kappa = 2/3$, the threshold in
\eqref{fr:eq:tstar} is $2 \times 0.335/0.125 = 5.36$, so
$T^{\circ *} = 6$, the number of waves in Section~\ref{P-sec:measurement}
of the paper. The exact value is $T^* = 8$, because $f_8/8 = 0.166$ is the
first value of $f_T/T$ at or below $0.186$. Carrying the bootstrap draws
through the two thresholds gives 95 percent intervals of $[5, 7]$ waves
under the approximation and $[6, 10]$ under the exact formula. Six biennial waves span a decade of panel coverage and eight span fourteen years, so the two-thirds benchmark that \citet{Mazumder2005} uses for the United States is reached in this population after a decade or more of biennial coverage, as the paper says.

\subsection{Rank-rank slopes}

Section~\ref{P-sec:measurement} of the paper reports by simulation that a
single-wave design recovers 0.52 of the true rank-rank slope, roughly
$\sqrt{\lambda_1}$, and 0.36 once the child's rank is itself measured from
three waves. Table~\ref{oatab:rank} gives the simulated
factors for one to seven waves. The following result gives the closed
form behind the simulation and explains the square root.

\begin{proposition}[Rank-rank attenuation under normality]
\label{fr:prop:rank}
Let permanent parental income $\alpha$ and the child outcome $c$ be
jointly normal with correlation $r \in (0,1]$, and let the measured
parental income be $m_T = \alpha + \bar v_T$ with $\bar v_T$ normal,
independent of $(\alpha, c)$, and reliability
$\Var(\alpha)/\Var(m_T) = \lambda_T$. Then the population rank-rank slope of
$c$ on $m_T$, relative to the slope of $c$ on $\alpha$, is
\begin{equation}
R_T \;=\; \frac{\arcsin\!\big(r\sqrt{\lambda_T}/2\big)}{\arcsin(r/2)},
\qquad \text{and} \qquad r\to 0 \;\Rightarrow\; R_T \to \sqrt{\lambda_T},
\label{fr:eq:rank}
\end{equation}
with $R_T \le \sqrt{\lambda_T}$ for every $r \in (0,1]$. If the child outcome
is itself measured as a $T_c$-wave average $c_{T_c} = c + \bar w_{T_c}$,
with $\bar w_{T_c}$ normal, independent of $(\alpha, c, \bar v_T)$, and
reliability $\Var(c)/\Var(c_{T_c}) = \lambda_{T_c}$, the factor becomes
\[
\frac{\arcsin\!\big(r\sqrt{\lambda_T\,\lambda_{T_c}}/2\big)}{\arcsin(r/2)},
\]
which tends to $\sqrt{\lambda_T\,\lambda_{T_c}}$ as $r \to 0$.
\end{proposition}

\begin{proof}\sloppy
Spearman's rank correlation for a bivariate normal pair with correlation
$\varrho$ is $(6/\pi)\arcsin(\varrho/2)$ \citep{Pearson1907, Kruskal1958}.
Population ranks are uniform on $[0,1]$ on both sides, so the rank-rank
regression slope equals the rank correlation. The pair $(m_T, c)$ is
bivariate normal, and since
\[
\Cov(m_T, c) = \Cov(\alpha, c) \quad\text{and}\quad \Var(m_T) = \Var(\alpha)/\lambda_T,
\]
its correlation is $\Corr(m_T, c) = r\sqrt{\lambda_T}$; taking the ratio of
the two Spearman correlations gives \eqref{fr:eq:rank}. As $r \to 0$,
$\arcsin(x) \sim x$ gives the square root. The inequality follows because
$\arcsin(x)/x$ is increasing on $(0,1]$, so
\[
\frac{\arcsin(r\sqrt{\lambda_T}/2)}{r\sqrt{\lambda_T}/2} \;\le\; \frac{\arcsin(r/2)}{r/2},
\]
which rearranges to $R_T \le \sqrt{\lambda_T}$. The two-sided statement replaces
the correlation $r\sqrt{\lambda_T}$ with
\[
\Corr(m_T, c_{T_c}) = r\sqrt{\lambda_T}\sqrt{\lambda_{T_c}}.
\]
\end{proof}

\begin{table}[!htb]
\centering
\small
\caption{Rank-rank attenuation: closed form against simulation.}
\label{fr:tab:rank}
\begin{tabular}{@{}lccccccc@{}}
\toprule
Waves of parental income $T$ & 1 & 2 & 3 & 4 & 5 & 6 & 7 \\
\midrule
$\lambda_T$ (exact) & 0.271 & 0.391 & 0.476 & 0.540 & 0.590 & 0.631 & 0.664 \\
$\sqrt{\lambda_T}$ & 0.521 & 0.625 & 0.690 & 0.735 & 0.768 & 0.794 & 0.815 \\
Closed form \eqref{fr:eq:rank}, $r = 0.30$ & 0.519 & 0.624 & 0.688 & 0.734 & 0.767 & 0.793 & 0.814 \\
Simulation & 0.517 & 0.620 & 0.685 & 0.734 & 0.770 & 0.794 & 0.814 \\
\midrule
Closed form, child measured with 3 waves & 0.358 & 0.430 & 0.474 & 0.505 & 0.529 & 0.546 & 0.561 \\
Simulation, child measured with 3 waves & 0.358 & 0.429 & 0.473 & 0.506 & 0.532 & 0.549 & 0.561 \\
\bottomrule
\end{tabular}
\par\smallskip
\begin{minipage}{0.95\textwidth}
\footnotesize\emph{Notes:} The simulation rows are the factors of Online
Appendix Table~\ref{oatab:rank}, printed to three decimals (400,000
Gaussian draws, $r = 0.30$). The
closed-form rows evaluate \eqref{fr:eq:rank} at the $\lambda_T$ of the first
row, which is the value the simulation uses. That value is computed from the unrounded pooled estimates and is the exact column of Table~\ref{fr:tab:etaboot}. The largest discrepancy between closed form and simulation is
0.004.
\end{minipage}
\end{table}

Table~\ref{fr:tab:rank} shows that the closed form reproduces the
simulated factors to within 0.004. Under normality the square-root rule holds up to a correction, $R_T \le \sqrt{\lambda_T}$, of at most 0.002 at $r = 0.30$. Two consequences follow without
further simulation. First, ranks attenuate less than levels, since
$\sqrt{\lambda_T} > \lambda_T$, but a single-wave rank design at
$\lambda_1 = 0.27$ still recovers only half the slope. Second, the factor
with both sides measured with error is a product, so transitory variation
in the child's income lowers rank slopes as much as transitory variation
in the parents' income does, whereas Section~\ref{fr:sec:side} shows
that it leaves slopes in levels unchanged.

\subsection{Delta-method standard errors}

The bootstrap in Table~\ref{fr:tab:etaboot} can be replaced by a
first-order expansion when only the parameter covariance is available.
Writing $\theta = (\sigma^2_\eta, \sigma^2_\varepsilon, \rho)$,
$D_T = \sigma^2_\eta + \sigma^2_v f_T/T$, and
$g_T(\rho) = f_T(\rho)/(1-\rho^2)$, so that
$\lambda_T = \sigma^2_\eta/(\sigma^2_\eta + \sigma^2_\varepsilon g_T/T)$,
\begin{equation}
\frac{\partial \lambda_T}{\partial \sigma^2_\eta} = \frac{\sigma^2_\varepsilon g_T/T}{D_T^2}, \qquad
\frac{\partial \lambda_T}{\partial \sigma^2_\varepsilon} = -\frac{\sigma^2_\eta\, g_T/T}{D_T^2}, \qquad
\frac{\partial \lambda_T}{\partial \rho} = -\frac{\sigma^2_\eta\sigma^2_\varepsilon\, g_T'(\rho)/T}{D_T^2},
\label{fr:eq:delta}
\end{equation}
with $g_T'(\rho) = \big[f_T'(\rho)(1-\rho^2) + 2\rho f_T(\rho)\big]/(1-\rho^2)^2$
and $f_T'(\rho) = 2\sum_{k=1}^{T-1}k(1-k/T)\rho^{k-1}$. The delta-method
variance is $\nabla\lambda_T' \,\hat V_\theta\, \nabla\lambda_T$. The
minimum-distance estimator imposes $\sigma^2_\eta \ge 0$, and the estimate
for advantaged children lies close to that boundary
(Table~\ref{P-tab:md} and Figure~\ref{P-fig:ridge} of the paper). Near that boundary neither the expansion nor the percentile bootstrap applies (Remark~\ref{fr:prop:boundary}), so Table~\ref{fr:tab:etaboot} reports the pooled sample only.

\section{The minimum-distance estimator}
\label{fr:sec:md}

\subsection{Moments, parameters, and the estimator}

On the five waves of 2014 to 2022 the estimator uses the $K = 15$ pairwise
autocovariances $\hat\omega_{ts}$, $t \le s$, where $t$ and $s$ index
waves. Each is the sample mean of $\tilde y_{it}\tilde y_{is}$ over the
$n_{ts}$ children observed in both waves, after demeaning within year
(and within group for the fits by group). The model \eqref{fr:eq:omega}
implies $\omega_{ts} = m_{s-t}(\theta)$ with
$m_k(\theta) = \sigma^2_\eta + \rho^k\sigma^2_\varepsilon/(1-\rho^2)$
and $\theta = (\sigma^2_\eta, \sigma^2_\varepsilon, \rho)$, so the model
has $p = 3$ parameters and imposes $K - p = 12$ restrictions. The
estimator weights each moment by its cell count $n_{ts}$ and solves
\begin{equation}
\hat\theta = \arg\min_{\theta \in \Theta}\; Q_N(\theta), \qquad
Q_N(\theta) = \sum_{t \le s} \frac{n_{ts}}{N}\big(\hat\omega_{ts} - m_{s-t}(\theta)\big)^2,
\qquad J = N\,Q_N(\hat\theta),
\label{fr:eq:md}
\end{equation}
with $N = \sum_{t\le s} n_{ts}$ and
$\Theta = \{\sigma^2_\eta \ge 0,\ \sigma^2_\varepsilon > 0,\ |\rho| < 1\}$.
This is minimum distance with the fixed diagonal weight matrix
$W = \mathrm{diag}(n_{ts}/N)$ \citep{Chamberlain1984, NeweyMcFadden1994}.

\begin{assumption}[Sampling]
\label{fr:ass:sampling}
Children are an i.i.d.\ draw; the pattern of observed waves is independent of $(\eta_i, \{\varepsilon_{it}\})$; and fourth moments of $\tilde y$ are
finite.
\end{assumption}

Under Assumption~\ref{fr:ass:sampling} and the model,
$\sqrt{n}(\hat\omega - \omega(\theta_0)) \Rightarrow N(0, \Omega)$
for the vector of pairwise moments, where $n$ is the number of children,
\[
\Omega_{(ts),(t's')} \;=\; \frac{\pi_{ts,t's'}}{\pi_{ts}\,\pi_{t's'}}\,
\Big(\E[\tilde y_t\tilde y_s\tilde y_{t'}\tilde y_{s'}] - \omega_{ts}\,\omega_{t's'}\Big),
\]
$\pi_{ts}$ is the probability that a child is observed in both waves $t$
and $s$, and $\pi_{ts,t's'}$ is the probability that a child is observed
in all of the waves $t$, $s$, $t'$, and $s'$. The independence part of the
assumption is what lets pairwise deletion estimate population
autocovariances. It is a missing-at-random condition on wave participation. The screens that define the sample pass 45.5 percent of advantaged and 42.0 percent of less advantaged children (Table~\ref{oatab:screens}); the assumption concerns which waves a child in the sample is observed in, which the screens do not test.

\subsection{Identification}

\begin{proposition}[Identification from three lags]
\label{fr:prop:ident}
Suppose the autocovariances at three consecutive lags $k_1$,
$k_2 = k_1 + 1$, and $k_3 = k_1 + 2$ are known. On $\Theta$ with
$\sigma^2_\varepsilon > 0$, and with $\rho \ne 0$ if $k_1 \ge 1$, the
parameter $\theta$ is identified:
\begin{equation}
\begin{aligned}
\rho &= \frac{\omega_{k_2} - \omega_{k_3}}{\omega_{k_1} - \omega_{k_2}}, &
\sigma^2_v &= \frac{\omega_{k_1} - \omega_{k_2}}{\rho^{k_1}(1 - \rho)}, \\
\sigma^2_\eta &= \omega_{k_1} - \rho^{k_1}\sigma^2_v, &
\sigma^2_\varepsilon &= \sigma^2_v(1-\rho^2).
\end{aligned}
\label{fr:eq:ident}
\end{equation}
With the variance among the three moments ($k_1 = 0$), identification
fails only when $\sigma^2_\varepsilon = 0$, in which case all
autocovariances equal $\sigma^2_\eta$ and $\rho$ is undetermined. It is
lost in the limit $\rho \to 1$, where the autocovariance function becomes
flat and the permanent variance cannot be separated from a transitory
component that never decays.
\end{proposition}

\begin{proof}
From \eqref{fr:eq:omega},
$\omega_{k_1} - \omega_{k_2} = \rho^{k_1}(1-\rho)\sigma^2_v$ and
$\omega_{k_2} - \omega_{k_3} = \rho^{k_1+1}(1-\rho)\sigma^2_v$. The
denominator $\omega_{k_1} - \omega_{k_2}$ is nonzero when
$\sigma^2_v > 0$, $\rho \ne 1$, and either $k_1 = 0$ or $\rho \ne 0$, and
the ratio is then $\rho$. The remaining expressions give $\sigma^2_v$,
$\sigma^2_\eta$, and $\sigma^2_\varepsilon$ in turn.
\end{proof}

The Jacobian $G = \partial\omega(\theta)/\partial\theta'$ has rows
$\big(1,\ \rho^k/(1-\rho^2),\ \sigma^2_\varepsilon\,\partial_\rho[\rho^k/(1-\rho^2)]\big)$.
The determinant of its three rows at lags 0, 1, and 2 is
$\sigma^2_\varepsilon/(1+\rho)^2$, so $G$ has full column rank whenever
$\sigma^2_\varepsilon > 0$ and $|\rho| < 1$. The rank condition holds at
every estimate in Table~\ref{P-tab:md} of the paper, and the ridge of
Section~\ref{P-sec:md} of the paper is a matter of the conditioning of $G$
and not of its rank. The next result quantifies the conditioning.

\subsection{The identification ridge}

Section~\ref{P-sec:md} of the paper reports that $\sigma^2_\eta$ and
$\rho$ trade off along a ridge of the criterion, and that the bootstrap
draws for advantaged children spread along it
(Figure~\ref{P-fig:ridge} of the paper). The following result describes
that curve and shows why it flattens as persistence rises.

\begin{proposition}[Leverage of the higher-order autocorrelations along the level set of the first-order autocorrelation]
\label{fr:prop:ridge}
Write the model in autocorrelation form,
$r_k = \omega_k/\omega_0 = s + (1-s)\rho^k$. Fix $r_1$ and move along the
set of $(s, \rho)$ pairs consistent with it,
$\rho(s) = (r_1 - s)/(1-s)$. Then the $k$-th autocorrelation changes
along this set at the rate
\begin{equation}
\frac{d r_k}{d s}\Big|_{r_1\ \mathrm{fixed}} \;=\; h_k(\rho) \;\equiv\; 1 - \rho^{k} - k\rho^{k-1}(1-\rho)
\;=\; (1-\rho)^2\sum_{j=0}^{k-2}(j+1)\rho^{j},
\qquad k \ge 2,
\label{fr:eq:hk}
\end{equation}
so that $h_2 = (1-\rho)^2$, $h_3 = (1-\rho)^2(1+2\rho)$, and every $h_k$
vanishes quadratically as $\rho \to 1$.
\end{proposition}

\begin{proof}
Differentiate $r_k(s) = s + (1-s)\rho(s)^k$ using
$\rho'(s) = -(1-r_1)/(1-s)^2 = -(1-\rho)/(1-s)$:
$dr_k/ds = 1 - \rho^k + (1-s)k\rho^{k-1}\rho'(s) = 1 - \rho^k - k\rho^{k-1}(1-\rho)$.
For the factorization, write $1 - \rho^k = (1-\rho)\sum_{j=0}^{k-1}\rho^j$
and $k\rho^{k-1}(1-\rho) = (1-\rho)\sum_{j=0}^{k-1}\rho^{k-1}$, so
$h_k = (1-\rho)\sum_{j=0}^{k-1}(\rho^j - \rho^{k-1}) = (1-\rho)\sum_{j=0}^{k-2}\rho^j(1-\rho^{k-1-j})$,
and expanding $1 - \rho^{k-1-j} = (1-\rho)\sum_{i=0}^{k-2-j}\rho^i$ and
collecting powers of $\rho$ gives the stated sum.
\end{proof}

The first-order autocorrelation fixes one combination of the permanent
share and the persistence parameter. The autocorrelations of second and
higher order separate the two, and they do so with leverage
$h_k(\rho)$. At the estimate for less advantaged children, $\hat\rho = 0.12$, $h_2 = 0.77$: moving the permanent share by one unit along the level set
of $r_1$ moves $r_2$ by 0.77, so the second-order moment determines the
split tightly. At the estimate for advantaged children,
$\hat\rho = 0.33$, $h_2 = 0.45$ and $h_3 = 0.75$. The leverage of the second-order moment is little more than half as large, and it is applied to moments from
420 children instead of 2,103, whose standard errors are larger by a
factor of about 2.2 on sample size alone. Weaker leverage and noisier
moments together produce the elongated cloud of bootstrap draws in
Figure~\ref{P-fig:ridge} of the paper. Proposition~\ref{fr:prop:ridge}
shows that the elongation is a property of the model at high $\rho$ and
not of the estimator: any method that identifies the permanent share
from the decay of the autocovariance function faces the same
$(1-\rho)^2$ factor.

\subsection{Asymptotic distribution and the choice of weights}

\begin{proposition}[Asymptotics of the estimator with fixed weights]
\label{fr:prop:asy}
Under Assumption~\ref{fr:ass:sampling}, if $\theta_0$ is the unique minimizer of the population criterion on $\Theta$ (Proposition~\ref{fr:prop:ident}), $G$ has full column rank (shown above), $\theta_0$ is interior to $\Theta$, $\Omega$ is positive definite, and the weight matrix converges to a positive definite $W$,
\begin{equation}
\sqrt{n}(\hat\theta - \theta_0) \Rightarrow N\!\big(0,\ (G'WG)^{-1}G'W\Omega WG(G'WG)^{-1}\big),
\label{fr:eq:asy}
\end{equation}
and the criterion $J = N Q_N(\hat\theta)$ converges in distribution to
$\sum_{j=1}^{K-p}\nu_j\chi^2_{1,j}$, a weighted sum of independent
$\chi^2_1$ variables whose weights $\nu_j$ are the nonzero
eigenvalues of $\Omega^{1/2}\big(W - WG(G'WG)^{-1}G'W\big)\Omega^{1/2}$
multiplied by the limit of $N/n$. For the optimally weighted criterion
$n\,(\hat\omega - \omega(\theta))'\hat\Omega^{-1}(\hat\omega - \omega(\theta))$
the weights all equal one and the limit is $\chi^2_{K-p}$
\citep{Hansen1982}. For the fixed weights of \eqref{fr:eq:md} the weights
$\nu_j$ depend on $\Omega$, and the $\chi^2_{K-p}$ reference
distribution does not apply.
\end{proposition}

\begin{proof}
The argument is the standard one for minimum distance \citep{NeweyMcFadden1994}; consistency follows from their Theorem 2.1 on a compact subset of $\Theta$ that contains $\theta_0$ in its interior, where the population criterion is continuous and has its unique minimum at $\theta_0$. The first-order condition
$G'W(\hat\omega - \omega(\hat\theta)) = 0$ and a mean-value expansion
give \eqref{fr:eq:asy}. For the criterion, the minimized residual is
$(I - G(G'WG)^{-1}G'W)(\hat\omega - \omega(\theta_0)) + o_p(n^{-1/2})$,
a linear map of an asymptotically normal vector, so
$J = (N/n)\,z'\,\Omega^{1/2}\big(W - WG(G'WG)^{-1}G'W\big)\Omega^{1/2}z + o_p(1)$
with $z$ asymptotically standard normal. A quadratic form $z'Az$ in a
standard normal vector with $A$ symmetric is distributed as
$\sum_j \nu_j\chi^2_{1,j}$ with $\nu_j$ the eigenvalues of $A$,
and the matrix here has rank $K - p$. With $W = \Omega^{-1}$ it is a
projection, whose nonzero eigenvalues equal one.
\end{proof}

Section~\ref{P-sec:md} of the paper reports $J$ under fixed cell-count
weights as a measure of fit and tests the model with a parametric
bootstrap of the inverse-variance-weighted criterion. Proposition~\ref{fr:prop:asy} is the reason for both choices. The
weights $\nu_j$ depend on $\Omega$, and with fifteen pairwise moments
from an unbalanced panel of 2,523 children the estimate $\hat\Omega$ is
noisy. \citet{AltonjiSegal1996} show that the optimally weighted
estimator is biased in small samples because $\hat\Omega$ is correlated
with the sample moments it weights. Fixed weights give up efficiency to
avoid that bias. Inverse-variance reweighting of the
moments moves the pooled permanent share from 0.271 to 0.264 (Section~\ref{P-sec:md} of the paper), so the estimate is not
sensitive to the choice between the two sets of weights.

\subsection{Decomposing the criterion}

The twelve overidentifying restrictions are of two kinds. Ten of them
say that autocovariances at the same lag but different calendar dates
are equal (stationarity). Two of them say that the five lag-specific
autocovariances lie on the curve $\sigma^2_\eta + \rho^k\sigma^2_v$
(shape). With a diagonal weight matrix the criterion splits exactly
into the two.

\begin{proposition}[Stationarity and shape components of $J$]
\label{fr:prop:Jsplit}
Let $\mathcal P_k$ be the set of wave pairs at lag $k$,
$W_k = \sum_{(t,s)\in\mathcal P_k}n_{ts}$, and
$\bar\omega_k = W_k^{-1}\sum_{(t,s)\in\mathcal P_k}n_{ts}\hat\omega_{ts}$
the mean of the moments at lag $k$ weighted by cell counts. Then for
every $\theta$,
\begin{equation}
N Q_N(\theta) \;=\; \underbrace{\sum_{k}\sum_{(t,s)\in\mathcal P_k} n_{ts}\big(\hat\omega_{ts} - \bar\omega_k\big)^2}_{J_{\mathrm{stat}}}
\;+\; \underbrace{\sum_{k} W_k\big(\bar\omega_k - m_k(\theta)\big)^2}_{J_{\mathrm{shape}}(\theta)} ,
\label{fr:eq:Jsplit}
\end{equation}
where $J_{\mathrm{stat}}$ does not depend on $\theta$. Consequently the
estimator on the fifteen pairwise moments coincides with the estimator
that fits the five lag means with weights $W_k$, and
$J = J_{\mathrm{stat}} + J_{\mathrm{shape}}(\hat\theta)$ with
$K - L = 10$ and $L - p = 2$ degrees of freedom attached to the two
components, $L = 5$ being the number of distinct lags.
\end{proposition}

\begin{proof}
Within lag $k$, write
$\hat\omega_{ts} - m_k = (\hat\omega_{ts} - \bar\omega_k) + (\bar\omega_k - m_k)$
and expand the weighted square; the cross term is
$2(\bar\omega_k - m_k)\sum_{\mathcal P_k}n_{ts}(\hat\omega_{ts} - \bar\omega_k) = 0$
by the definition of the weighted mean. Summing over $k$ gives
\eqref{fr:eq:Jsplit}. Since the first term is free of $\theta$, minimizing
$Q_N$ is the same as minimizing $J_{\mathrm{shape}}$.
\end{proof}

\begin{table}[!htb]
\centering
\small
\setlength{\tabcolsep}{4pt}
\caption{Decomposition of the minimum-distance criterion.}
\label{fr:tab:Jsplit}
\begin{tabular}{@{}lccccc@{}}
\toprule
& $J$ & $J_{\mathrm{stat}}$ & $J_{\mathrm{shape}}$ & Stationarity & Lag means \\
& (12 df) & (10 df) & (2 df) & share of $J$ & $\bar\omega_0,\ldots,\bar\omega_4$ \\
\midrule
Pooled & 29.31 & 28.85 & 0.46 & 0.984 & 0.459, 0.177, 0.141, 0.123, 0.110 \\
Advantaged children & 2.66 & 2.51 & 0.14 & 0.946 & 0.438, 0.198, 0.128, 0.105, 0.066 \\
Less advantaged children & 29.76 & 29.51 & 0.25 & 0.992 & 0.460, 0.169, 0.141, 0.124, 0.120 \\
\bottomrule
\end{tabular}
\par\smallskip
\begin{minipage}{0.95\textwidth}
\footnotesize\emph{Notes:} Equation \eqref{fr:eq:Jsplit} at the
minimum-distance estimates of Table~\ref{P-tab:md} of the paper. $J$
reproduces the paper's values (29.3, 2.7, 29.8). Degrees of freedom are
nominal: by Proposition~\ref{fr:prop:asy} neither component is $\chi^2$
under fixed weights, and the paper reports the total as a measure of fit only.
\end{minipage}
\end{table}

Table~\ref{fr:tab:Jsplit} reports the split. The stationarity component
is 98.4 percent of the criterion in the pooled sample and 99.2 percent for less advantaged children. Autocovariances at the same lag differ across
calendar dates, most visibly the variances, which are 0.58 in 2016 against 0.41 to 0.49 in the other waves (pooled row of
Table~\ref{oatab:autocov}). The shape component is 0.46 in
the pooled sample: a permanent effect plus an AR(1) component reproduces
the five lag means closely. The decomposition locates the
misfit that the paper reports for the pooled sample and for less advantaged children. It lies in the departures of same-lag moments from equality across calendar dates, the sense in which the paper says that stationarity holds only approximately, and not in the shape of the curve. The estimates depend on
the data only through the five lag means
(Proposition~\ref{fr:prop:Jsplit}), so the variation across calendar
dates does not enter them; when stationarity fails they describe the
autocovariance function averaged over waves. For advantaged children the
criterion is 2.66 and has the same composition, with 94.6 percent in the stationarity component. With 420 children, departures from stationarity
of a given size add less to a criterion whose weights are cell counts.

\subsection{Inference at the boundary}

\begin{remark}[Bootstrap inconsistency at the boundary]
\label{fr:prop:boundary}
Let $\hat\theta$ be the constrained estimator on $\Theta$ with
$\sigma^2_\eta \ge 0$. If $\sigma^2_{\eta,0} = 0$, or if
$\sqrt{n}\,\sigma^2_{\eta,0}$ is bounded, the limit distribution of
$\sqrt{n}(\hat\sigma^2_\eta - \sigma^2_{\eta,0})$ is that of a normal
variable censored from below, and the nonparametric bootstrap
distribution of $\hat\sigma^2_\eta$ does not converge to it
\citep{Andrews2000}. Percentile intervals for $\sigma^2_\eta$, and for
functions of it such as the permanent share and $\lambda_T$, are then not
asymptotically valid. \citet{Andrews2000} establishes the result for the mean of a normal
sample restricted to be nonnegative, and the argument carries over: when
the true parameter is on the boundary or within $O(n^{-1/2})$ of it, the
bootstrap does not replicate the mass that the constrained estimator
places at the boundary, because the bootstrap is centered at
$\hat\theta$ and not at $\theta_0$.
\end{remark}

The row of draws for advantaged children at a permanent share of zero in
Figure~\ref{P-fig:ridge} of the paper is this boundary. With
$\hat\sigma^2_\eta = 0.083$ and a bootstrap standard error of 0.028, the
estimate for advantaged children is about three standard errors from the
boundary, and the bootstrap draws reach it; this is the situation of
Remark~\ref{fr:prop:boundary}. The bootstrap standard
errors that Table~\ref{P-tab:md} of the paper reports for that group
carry this qualification, which is the reason
Section~\ref{P-sec:md} of the paper gives for not reading the ratio of a
group gap to its standard error at face value. The paper treats the
fitted split for advantaged children as suggestive and rests the
comparison of groups on
the fifteen raw moments (Figure~\ref{P-fig:autocov} of the paper, Table~\ref{oatab:autocov}, and the
joint test of Section~\ref{oa:jointtest}). Their
estimators are unconstrained sample means, to which the bootstrap applies
without qualification.

\subsection{Observed and transitory persistence}

Under \eqref{fr:eq:model}, the population first-order autocorrelation of
$\tilde y$ is
\begin{equation}
\rho_{\mathrm{obs}} \;=\; \Corr(\tilde y_{it}, \tilde y_{i,t+1}) \;=\; s + (1-s)\rho,
\label{fr:eq:rhoobs}
\end{equation}
exactly: the ratio of $\Cov(\tilde y_t, \tilde y_{t+1}) = \sigma^2_\eta + \rho\sigma^2_v$ to
$\Var(\tilde y_t) = \sigma^2_\eta + \sigma^2_v$, with $s$ substituted, is \eqref{fr:eq:rhoobs}.
It is the probability limit of the pooled least squares coefficient of
$\tilde y_{it}$ on $\tilde y_{i,t-1}$ without individual effects, which converges
to the ratio of these two moments, and it exceeds $\rho$ whenever $s > 0$.

Equation~\eqref{P-eq:rhoobs} of the paper states \eqref{fr:eq:rhoobs} under stationarity. In the population it is exact; applied to a pooled regression on an unbalanced panel whose waves have unequal variances (Table~\ref{oatab:autocov}) it holds approximately. At the
pooled estimates, $\rho_{\mathrm{obs}} = 0.27 + 0.73 \times 0.16 = 0.39$,
the value in Section~\ref{P-sec:abgmm} of the paper.

The Arellano--Bond estimator of Section~\ref{P-sec:abgmm} of the paper
targets $\rho$ and not $\rho_{\mathrm{obs}}$. Rewriting
\eqref{fr:eq:model} as
$\tilde y_{it} = \rho\tilde y_{i,t-1} + (1-\rho)\eta_i + \varepsilon_{it}$
shows that the autoregression in equation~\eqref{P-eq:ab} of the paper
has individual effect $\alpha_i = (1-\rho)\eta_i$ and idiosyncratic error
$\varepsilon_{it}$, which is serially uncorrelated, so lagged levels dated
$t-2$ and earlier are valid instruments for the first-differenced
equation \citep{ArellanoBond1991}. The two estimators share a target, as
the paper says, so the gap between the minimum-distance estimate and the
Arellano--Bond estimate is not a difference of definitions; the paper notes that the minimum-distance estimate for less advantaged children, 0.12, lies inside the 95 percent interval of the Arellano--Bond estimate for that group, $[-0.01, 0.13]$.

\section{The accounting identity for the premium}
\label{fr:sec:identity}

Section~\ref{P-sec:mediation} of the paper decomposes the premium through identity~\eqref{P-eq:hp-intro} of the paper, $\tau = \beta^{H}\,\Delta M + \Direct$, where $\tau$ is the premium with parental income held fixed, $\Delta M$ is the gap in college completion between advantaged and less advantaged children, and $\beta^{H}$ is the return to college for advantaged children. The paper evaluates the identity at the least squares returns and at the instrumented return to college. The following result states the identity in observed conditional means and shows what the direct term collects and what the remainder contains when the identity is evaluated at a return other than $\beta^{H}$.
Expectations for less advantaged children are taken under the entropy-balancing
weights of Section~\ref{P-sec:ebal} of the paper, which match the first
two moments of parental income, gender, age, urban status, five-year birth cohort, and province (the fifteen largest separately) to those of advantaged children.

\begin{proposition}[Content of the direct term]
\label{fr:prop:identity}
Let $H \in \{0,1\}$ indicate advantaged children and $M \in \{0,1\}$ college
completion. Define $p_h = \Pr(M = 1 \mid H = h)$,
$a_h = \E[Y \mid H = h, M = 0]$, the mean log earnings of the children of
group $h$ without a degree, and
$r_h = \E[Y \mid H = h, M = 1] - \E[Y \mid H = h, M = 0]$, the earnings
difference between the graduates of group $h$ and its children without a
degree. Then, with $\tau = \E[Y \mid H = 1] - \E[Y \mid H = 0]$,
$\Delta M = p_1 - p_0$, and $\beta^{H} = r_1$ as in
Section~\ref{P-sec:mediation} of the paper:
\begin{enumerate}
\item[(i)] The identity holds with
\begin{equation}
\tau \;=\; r_1\,\Delta M \;+\; \Direct, \qquad
\Direct \;=\; (a_1 - a_0) \;+\; p_0\,(r_1 - r_0)
\;=\; (a_1 + p_0\,r_1) - (a_0 + p_0\,r_0) .
\label{fr:eq:identity}
\end{equation}
The direct term is the sum of the gap between the groups in earnings
without college and a term equal to the college completion rate of less advantaged children, $p_0$, times the difference in returns between the two groups' graduates. It equals $a_1 - a_0$ if and only if $r_1 = r_0$ or
$p_0 = 0$.
\item[(ii)] For any number $b$,
\begin{equation}
\tau - b\,\Delta M \;=\; \Direct \;+\; (r_1 - b)\,\Delta M .
\label{fr:eq:remainder}
\end{equation}
\end{enumerate}
\end{proposition}

\begin{proof}
By the law of total expectation, $\E[Y \mid H = h] = (1 - p_h)\,a_h + p_h\,(a_h + r_h) = a_h + p_h r_h$. Subtract the expression for $h = 0$
from the expression for $h = 1$ and add and subtract $p_0 r_1$:
$\tau = (a_1 - a_0) + p_1 r_1 - p_0 r_0 = (a_1 - a_0) + (p_1 - p_0)r_1 + p_0(r_1 - r_0)$,
which is (i). Part (ii) adds and subtracts $b\,\Delta M$.
\end{proof}

The identity is an accounting identity in observed conditional means and
requires no assumption. Written in potential outcomes, with
$Y = Y(0) + M\,(Y(1) - Y(0))$, the same algebra holds with
$a_h = \E[Y(0) \mid H = h]$ and $r_h = \E[Y(1) - Y(0) \mid H = h, M = 1]$;
the observed $r_h$ then equals that return plus the selection term
$\E[Y(0) \mid H = h, M = 1] - \E[Y(0) \mid H = h, M = 0]$, which is why the
paper reads the identity at a least squares return as accounting and not as
a causal statement. The last expression in \eqref{fr:eq:identity} gives the
direct term the reading of Section~\ref{P-sec:mediation} of the paper, the
earnings advantage that would remain if the college gap were closed, in a
specific form: the completion rate of advantaged children is lowered from
$p_1$ to $p_0$, and the graduates removed have mean return $r_1$. The
components $\tau = 0.072$ and $\Delta M = 0.122$ are differences between
the two groups after entropy balancing, and the paper reads the premium
as an association. Inside the balanced comparison the identity is exact:
graduates out-earn children without a degree by $r_1 = 0.541$ among
advantaged children and by $r_0 = 0.517$ among reweighted less advantaged
children, with $p_1 = 0.3355$ and $p_0 = 0.2130$, so
$r_1\,\Delta M = 0.066$, $p_0\,(r_1 - r_0) = 0.005$, and
$a_1 - a_0 = 0.001$ (Section~\ref{P-sec:mediation} of the paper). A causal
return takes the place of $r_1$ when the identity is evaluated at the
instrument. The instrument gives an estimate of it for the advantaged
children the expansion moved into college, 1.14, with a 90 percent
Anderson--Rubin set that excludes every return between $-3.20$ and 0.60
(Table~\ref{P-tab:groups} of the paper), and an estimate for all children,
0.68, with a 95 percent Anderson--Rubin set that starts at 0.50
(Table~\ref{P-tab:mature} of the paper). The paper evaluates the identity
over this range of returns.

At the least squares return, 0.449, the remainder is
$0.072 - 0.449 \times 0.122 = 0.017$
(Table~\ref{P-tab:bounds} of the paper); at the instrumented returns
the college channel, 0.084 for all children and 0.140 for advantaged children,
exceeds the premium. By \eqref{fr:eq:remainder} this
remainder equals $\Direct$ if the least squares return equals $r_1$. If
the least squares return exceeds $r_1$, the remainder understates
$\Direct$ by 0.122 times the excess, and if it falls short of $r_1$ the
remainder overstates $\Direct$ by 0.122 times the shortfall.
Table~\ref{P-tab:bounds} of the paper evaluates the left side of
\eqref{fr:eq:remainder} at alternative values of $b$. The second term of
$\Direct$ in \eqref{fr:eq:identity} vanishes when returns among graduates
are equal in the two groups. In the balanced comparison the two returns are 0.541 and 0.517, so the
second term is 0.005 and the direct term exceeds the gap between the groups
in earnings without a degree, $a_1 - a_0 = 0.001$, by that amount. When
returns are higher for advantaged graduates ($r_1 > r_0$), $\Direct$ is an
upper bound on $a_1 - a_0$.
